\documentclass{beamer}

\usepackage{amsmath, amsfonts, amssymb}

\usetheme{Antibes}

\usecolortheme{dolphin}

\usefonttheme{professionalfonts}

\setbeamertemplate{navigation symbols}{}

\AtBeginSection[]{
  \begin{frame}
    \vfill
    \centering
    \begin{beamercolorbox}[sep=8pt,center,shadow=true,rounded=true]{title}
      \usebeamerfont{title}\insertsectionhead\par%
    \end{beamercolorbox}
    \vfill
  \end{frame}
}

\title{Calculus 2}
\author{Bobby}
\institute{IIIS, Tsinghua University}
\date{\today}

\begin{document}

\begin{frame}
    \titlepage
\end{frame}

\begin{frame}{Contents}
    \tableofcontents[hideallsubsections]
\end{frame}

\begin{section}{Limits and Continuity of Multivariate Functions}

\begin{subsection}{Problem 1}
    \begin{frame}
        \begin{block}{Problem 1}
            Determine the necessary and sufficient conditions on the coefficients $a, b,$ and $c$ such that the quadratic form $f(x, y) = ax^2 + bxy + cy^2$ is of the same order as $g(x, y) = x^2 + y^2$ as $(x, y) \to (0, 0)$.    
        \end{block}
    \end{frame}

    \begin{frame}
        \begin{block}{Hint 1}
            To prove two functions are of the same order, you must show that the limit of their ratio is bounded between two positive constants $m$ and $M$ as the variables approach the origin.
        \end{block}

        \pause
        \vspace{2em}
        \begin{block}{Hint 2}
            Utilize the transformation to polar coordinates: $x = r \cos \theta$ and $y = r \sin \theta$. Observe how the radius $r$ cancels out in the ratio of the two quadratic forms.
        \end{block}

    \end{frame}

    \begin{frame}
        \begin{block}{Hint 3}
            Recall that a continuous function on a compact set (like $\theta \in [0, 2\pi]$) attains its maximum and minimum values. Use the condition $b^2 < 4ac$ to ensure the function never reaches zero.
        \end{block}
    \end{frame}
\end{subsection}

\begin{subsection}{Problem 2}
    \begin{frame}
        \begin{block}{Problem 2}
            Evaluate the following limit:
            \[
            \lim_{(x,y) \to (0,0)} \frac{x^5 + y^5}{x^3 + y^3}
            \]    
        \end{block}
    \end{frame}

    \begin{frame}
        \begin{block}{Hint 1}
            Start by factoring the numerator using the sum of fifth powers formula and the denominator using the sum of cubes formula. Look for a common algebraic factor, $(x+y)$, that can be canceled out to simplify the expression.
        \end{block}

        \pause
        \vspace{2em}
        \begin{block}{Hint 2}
            After simplifying, convert the remaining fractional expression into polar coordinates by letting $x = r \cos \theta$ and $y = r \sin \theta$. 
        \end{block}
    \end{frame}
\end{subsection}

\begin{subsection}{Problem 3}
    \begin{frame}
        \begin{block}{Problem 3}
            A set $K$ in $\mathbb{R}^m$ is called a convex set if for any $\mathbf{x}_1, \mathbf{x}_2 \in K$ and any $0 < t < 1$, we have $(1-t)\mathbf{x}_1 + t\mathbf{x}_2 \in K$. A function $f: K \to \mathbb{R}$ is called a convex function if for any $\mathbf{x}_1, \mathbf{x}_2 \in K$ and any $0 < t < 1$, we have
            $$
            f((1-t)\mathbf{x}_1 + t\mathbf{x}_2) \le (1-t)f(\mathbf{x}_1) + tf(\mathbf{x}_2).
            $$

            Prove that any convex function on an open convex set $K$ is locally Lipschitz, that is, for any $\mathbf{x}_0 \in K$, there exists a neighborhood $U$ of $\mathbf{x}_0$ and a constant $L > 0$ such that for any $\mathbf{x}, \mathbf{y} \in U$, we have
            $$
            |f(\mathbf{x}) - f(\mathbf{y})| \le L\|\mathbf{x} - \mathbf{y}\|.
            $$
        \end{block}
    \end{frame}

    \begin{frame}
        \begin{block}{Hint 1}
            Construct a polytope around $\mathbf{x}_0$ within $K$ to find a bound $M$ via vertex combinations. 
        \end{block}

        \pause
        \vspace{2em}
        \begin{block}{Hint 2}
            For any $\mathbf{x}, \mathbf{y} \in B(\mathbf{x}_0, r)$ with distance $d$, extend the ray from $\mathbf{y}$ through $\mathbf{x}$ to a point $\mathbf{z}$ at distance $r$ from $\mathbf{x}$. Verify $\mathbf{z} \in B(\mathbf{x}_0, 2r)$.
        \end{block}

        \pause
        \vspace{2em}
        \begin{block}{Hint 3}
            Express $\mathbf{x}$ as a convex combination: $\mathbf{x} = \frac{r}{r+d}\mathbf{y} + \frac{d}{r+d}\mathbf{z}$. Apply the convexity condition, subtract $f(\mathbf{y})$, and use bound $M$ to derive the Lipschitz constant $L = \frac{2M}{r}$.    
        \end{block}
    \end{frame}
\end{subsection}

\end{section}

\begin{section}{Differential of Multivariable Functions}

\begin{subsection}{Problem 1}
    \begin{frame}
        \begin{block}{Problem 1}
            Let $f(x,y)$ be a differentiable function on $\mathbb{R}^2 \setminus \{(0,0)\}$ and $k$ be a non-negative integer. Suppose that for any $(x,y) \neq (0,0)$, the following equation holds:
            $$
            x \frac{\partial f}{\partial x}(x,y) + y \frac{\partial f}{\partial y}(x,y) = k f(x,y)
            $$
            
            Prove that $f$ is a homogeneous function of degree $k$. That is, for any $(x,y) \neq (0,0)$ and any positive real number $t$, show that:
            $$
            f(tx, ty) = t^k f(x,y)
            $$
        \end{block}
    \end{frame}

    \begin{frame}
        \begin{block}{Hint 1}
            Fix an arbitrary point $(x,y) \neq (0,0)$. To connect $f(tx, ty)$ and $t^k f(x,y)$, construct an auxiliary function of a single variable $t > 0$:
            $$
            g(t) = t^{-k} f(tx, ty)
            $$
        
            Your ultimate goal is to prove that $g(t)$ is a constant function independent of $t$.
        \end{block}

        \pause
        \vspace{1em}
        \begin{block}{Hint 2}
            Calculate the derivative $g'(t)$ using the product rule and the multivariate chain rule. Then, evaluate the given condition (the partial differential equation) at the scaled point $X = tx$ and $Y = ty$. Substitute this relationship into your expression for $g'(t)$ to show that $g'(t) = 0$. Finally, evaluate $g(1)$ to finish the proof.    
        \end{block}
    \end{frame}
\end{subsection}
\end{section}

\end{document}